\section{Introduction}

\subsection{Clay Requirements}

The three-dimensional Navier--Stokes regularity problem asks whether a smooth
solution can lose its smoothness in finite time.  Fix smooth,
divergence-free, finite-energy initial data \(u_0\), fix \(\nu>0\), and let
\((u,p)\) be the corresponding maximal classical solution on \([0,T_*)\):
\[
  \partial_tu+(u\cdot\nabla)u+\nabla p=\nu\Delta u,
  \qquad \nabla\cdot u=0,
  \qquad u(0)=u_0.
\]
We want to prove that \(T_*\) cannot be finite.

The word \emph{maximal} has a precise meaning here.  The endpoint is
\[
  T_*
  :=
  \sup\bigl\{
    T>0:
    \text{the classical solution generated by \(u_0\) exists on \([0,T]\)}
  \bigr\}.
\]
For every \(T<T_*\), the same pair \((u,p)\) is smooth on \([0,T]\) and
satisfies the original equations there.  If \(T_*<\infty\), maximality says
that every classical pair on a later interval fails to agree with \((u,p)\)
throughout its preterminal history.  That failure of same-solution extension
is the definition of finite-time breakdown.  The pair remains smooth at every
\(t<T_*\); breakdown concerns its extension through the endpoint.

The Clay requirement concerns the full velocity--pressure pair.  In the
whole-space formulation, the initial velocity is smooth and rapidly decaying;
in the periodic formulation, the data and solution are periodic.  In either
case both \(u\) and \(p\) must remain smooth for all time, while the velocity
retains finite energy.  We choose the decaying pressure representative on
\(\mathbb R^3\) and the zero-mean representative on \(\mathbb T^3\).
Incompressibility then fixes the pressure response from the velocity; only
\(u\) is prescribed initially.  Taking the divergence of the momentum equation
and using \(\nabla\cdot u=0\) gives
\[
  -\Delta p=\partial_i\partial_j(u_i u_j).
\]
Thus the velocity field determines the simultaneous pressure response.  Before
normalization it is unique up to a time-dependent additive constant; the chosen
representative fixes that freedom.  A continuation criterion stated in terms
of \(u\) still includes pressure through this relation: the required velocity
regularity supplies the corresponding pressure regularity.

Local well-posedness turns non-extension into a quantity the proof can
estimate.  Fix \(s>5/2\) on the active spatial domain.  If
\[
  \sup_{0\le t<T_*}\|u(t)\|_{H^s}<\infty,
\]
then the same solution has enough regularity to reach the endpoint and restart
beyond it.  A finite maximal time must therefore force
\[
  \limsup_{t\uparrow T_*}\|u(t)\|_{H^s}=\infty
  \qquad\text{for every fixed }s>5/2.
\]
The continuation theorem makes this divergence a necessary consequence of
finite breakdown; the definition itself remains the failure of extension.  A
discontinuity or an unbounded velocity derivative is one possible
manifestation, and either one obstructs classical continuation.

The proof must now obtain continuation-strength control from the equation and
the initial datum.  The first available estimate is the energy identity; the
Gold Standard begins by determining exactly how far that estimate reaches.

\clearpage
\subsection{Gold Standard Proof}

The Gold Standard route starts with the smooth solution on \([0,T_*)\) and
tries to prove directly that the control needed for continuation stays finite
as \(t\) approaches \(T_*\).  The solution is already smooth throughout this
interval.  The proof must obtain bounds that remain uniform as the solution
approaches the endpoint.

The energy identity supplies the starting estimate.  Pair the velocity
equation with \(u\), integrate over space, and use incompressibility and
integration by parts.  This gives
\[
  \|u(t)\|_{L^2}^2
  +2\nu\int_0^t\|\nabla u(r)\|_{L^2}^2\,dr
  =\|u_0\|_{L^2}^2.
\]
The first term is the \(L^2\) size of the velocity at time \(t\).  The integral
records the viscous dissipation accumulated up to that time, and the right side
is fixed by the initial data.  The identity keeps the velocity bounded in
\(L^2\) and gives time-integrated \(L^2\) control of its first spatial
derivatives.

The proof then turns to higher derivatives of the same solution.  Here,
bootstrapping helps carry control from one derivative level to the next.  At
each step, the estimates already established are used to control the terms in
the differentiated equation.  When that estimate closes, its new bound becomes
available for the following step.  Sobolev \(H^s\) spaces organize the
\(L^2\) control of weak derivatives, and Sobolev embedding converts sufficiently
strong \(H^s\) control into the pointwise boundedness and continuity needed to
continue the classical solution.

The difficulty is closing this climb for every finite time.  The energy
identity supplies less control than the nonlinear transport term demands at
the next derivative levels.  A Gold Standard proof must find an a priori
estimate, derived from the initial data, the equation, and the viscosity, that
keeps the required Sobolev control finite.  Such an estimate would prevent the
failure condition in Section 1.1 and allow the same solution to continue past
\(T_*\), so a finite \(T_*\) could never be maximal.

The belief that viscosity can ``pay the whole bill'' is the belief that viscous
damping supplies enough control to absorb the nonlinear derivative growth at
every level needed by this argument.

\subsection{Silver Standard Proof}

Silver turns the problem around.  The Gold Standard tries to prove directly
that the equations keep the fluid smooth.  The Silver Standard begins with the
same conjecture and reads it from the other direction.  If the original
Navier--Stokes evolution always produces a smooth fluid, then a nonsmooth
outcome cannot still belong to that same evolution.

In shorthand, the idea is
\[
  \text{Navier--Stokes}\Longrightarrow\text{smooth}
  \qquad\Longleftrightarrow\qquad
  \text{not smooth}\Longrightarrow\text{not Navier--Stokes}.
\]
The two implications are contrapositives, so proving the second would prove the
first.  Silver does not take the first implication for granted.  It chooses the
second as its proof route.

Suppose the solution appears to lose smoothness at a finite time.  For this to
be a genuine counterexample, the nonsmooth object must still be the same fluid
that began from \(u_0\) and evolved with the fixed viscosity \(\nu\) under the
original equations.  Silver asks whether those two facts can coexist.  If every
possible loss of smoothness also forces the object to stop satisfying the
original problem, or to stop being the same fluid, then none of those losses is
a Navier--Stokes counterexample.

This is where Silver gains leverage.  Gold has to find the estimates that keep
the solution smooth and explain how the equations achieve that control.  Silver
can leave that mechanism open.  Its task is to rule out every nonsmooth outcome,
either one by one or through a single principle that covers them all.  Once that
exclusion is complete, the original fluid has no nonsmooth endpoint.

\subsection{Class Membership and Exiting the Class}

The idea of a Navier--Stokes class comes from putting Gold and Silver together.
Gold asks whether everything built into the original problem, acting together, guarantees smoothness.
If it does, Silver turns the question around: a nonsmooth fluid cannot still satisfy the whole problem.
The class contains the fluids that do satisfy it, and membership means that the fluid still meets every one of its requirements.

Membership begins with the fluid described by the original problem.
It has constant density, moves incompressibly, and has one fixed positive viscosity; its velocity and pressure continue to obey the Navier--Stokes equations.
These conditions appear together in the motion: the fluid can change shape without changing its volume, pressure gradients push and redirect the flow, and viscosity transfers momentum between neighboring parts moving at different speeds.
Class membership means that the fluid continues to behave this way as one complete system.

Couette flow gives a familiar picture.
Place the fluid between two parallel surfaces moving at different speeds.
Viscosity carries tangential momentum through the fluid, producing a continuous shear profile from one layer to the next.

A loss of membership could involve losing exactly that viscous behavior.
Now compare Navier--Stokes with Euler.
Euler keeps the constant-density incompressible motion, nonlinear transport, and pressure, while removing \(\nu\Delta u\).
A moving portion of fluid still keeps its volume, but the viscous transfer from one layer to the next is gone.

Imagine two parts of the fluid touching and sliding past one another like two wooden boards.
They do not pass through each other, so their normal motion can agree and incompressibility can remain intact.
Their tangential velocities can still differ, creating a jump across the surface where they meet.

With compatible pressure and normal flux, this slip can satisfy the weak Euler equations.
Euler has no viscous term to carry tangential momentum across the surface and smooth the jump.
For fixed viscosity, the same jump is seen by \(\nu\Delta u\), so the discontinuous sheet is not a smooth Navier--Stokes state.
The slip has kept the rules shared by both equations and lost the rule that distinguishes them.

This is the idea behind viscosity paying the whole bill for smoothness.
If viscosity is the entire reason Navier--Stokes stays smooth, then every nonsmooth scenario should cross this same line: as smoothness is lost, the fluid also leaves the Navier--Stokes class.
When the shared Euler rules survive, the scenario lands in Euler.
If one of those shared rules also fails, it leaves both classes.
If the fluid became nonsmooth while every Navier--Stokes requirement still held, it would be a genuine counterexample.

\subsection{Velocity Failures}

The flat slip is only the first shape suggested by this comparison.
Wrap the same slipping surface around into a loop.
Locally, each part of the loop repeats the wooden-board picture; globally, the mismatch is organized into a spinning vortex shape.
Euler leaves open the possibility that such a nonsmooth vortex can retain incompressibility, normal compatibility, and weak momentum balance.
Viscosity acts across the surface of the loop, so the same sharp separation is not a smooth fixed-viscosity flow.

The loop shows why the possible failures have to be generalized.
One local slip can appear as a flat sheet, a closed vortex, or another geometry without changing the basic defect in the velocity.
Other scenarios may make part of the velocity field cease to participate in the surrounding motion, or may drive the velocity beyond every finite bound.
The useful question is what happens to the velocity itself, independent of the particular shape carrying the failure.
This gives the three working pictures Dead, Jump, and Blown.

\paragraph{Dead.}
Dead describes a point or region that becomes inactive relative to the moving fluid around it.
A zero velocity alone can be an ordinary smooth stagnation point.
The stronger failure occurs when the point or region ceases to participate in the same evolving velocity field.

\paragraph{Jump.}
Jump describes two approaches to the same location that produce different velocities.
The flat slip and the wrapped vortex are two geometries of this same failure.
Across the break, the fluid no longer has one smooth velocity field.

\paragraph{Blown.}
Blown describes the velocity, or one of its derivatives, growing without a finite bound as the failure time is approached.
This is the familiar blow-up picture.

Dead, Jump, and Blown say what has happened to the velocity.
Class membership says which equations the fluid still obeys after the failure.
Because smoothness belongs to one complete field, each local picture has to be compared with the surrounding velocity and its derivatives.

\subsection{Coherent Field Participation}

The velocity at one point is only the bottom of the comparison.  Smoothness
also depends on every time and spatial derivative attached to that point.  We
collect them in the full velocity tower
\[
  \mathcal J^\infty u(t,x)
  :=
  \left\{
    \partial_t^j\partial_x^\alpha u(t,x):
    j\ge0,
    \ \alpha\in\mathbb N_0^3
  \right\}.
\]
Here \(j\) counts time derivatives and the multi-index \(\alpha\) records the
number of derivatives taken in each spatial direction.  The set contains the
velocity itself and every derivative above it.

For a point to be fully dead, every level of its tower would have to vanish
while the surrounding fluid remained active:
\[
  \partial_t^j\partial_x^\alpha u(t_*,x_*)=0
  \qquad\text{for every }j\ge0
  \text{ and }\alpha\in\mathbb N_0^3.
\]
This equation says much more than \(u(t_*,x_*)=0\): it removes every local rate
of change at the point.  Even an all-zero tower does not prove a failure by
itself, because a smooth field can be flat at one point.  The failure would be
the inability of that tower to remain part of one smooth field with the active
towers around it.

For Jump, write
\(D^{j,\alpha}u=\partial_t^j\partial_x^\alpha u\), and approach the same
point along two sequences \((t_n,x_n^+)\) and \((t_n,x_n^-)\).  A jump occurs
at one level of the tower if
\[
  \lim_{n\to\infty}D^{j,\alpha}u(t_n,x_n^+)
  \ne
  \lim_{n\to\infty}D^{j,\alpha}u(t_n,x_n^-).
\]
The two sequences reach the same point, but the same derivative has two
different limiting values.  That derivative cannot be continuous there.

For Blown, one level of the tower grows without a finite local bound.  At the
affected point this has the form
\[
  \limsup_{(t,x)\to(t_*,x_*)}
  \left|D^{j,\alpha}u(t,x)\right|=\infty
  \qquad\text{for some }j\text{ and }\alpha.
\]
The display says that the selected derivative becomes arbitrarily large along
points approaching \((t_*,x_*)\), so it cannot supply a finite value there.

These comparisons give the one-coherent-field condition.  The towers at
different points do not need to be equal.  They must vary compatibly as the
derivatives of one velocity field: every approach to the same point must give
the same finite value at each derivative level.  Dead tests whether a fully
vanishing tower can remain joined to the active field, Jump tests whether two
approaches remain compatible, and Blown tests whether the tower remains finite.
The failure becomes genuine when the proposed readings can no longer belong to
one smooth velocity field.

\subsection{VPI-Law}

One coherent field requires more than matching values and derivatives.  Those
derivatives must also obey the same evolution law.  Viscosity, pressure, and
incompressibility act together on the same velocity field.  We use VPI as
shorthand for this joined law.

Start with the momentum equation and the incompressibility condition.  Taking
the divergence of the momentum equation, and using \(\nabla\cdot u=0\), gives
the pressure equation.  Written together, the three relations are
\[
  \partial_tu+(u\cdot\nabla)u+\nabla p=\nu\Delta u,
  \qquad
  \nabla\cdot u=0,
  \qquad
  -\Delta p=\partial_i\partial_j(u_i u_j).
\]
The first relation evolves the velocity.  The second preserves
incompressibility.  The third determines the pressure response from the whole
velocity field.  The transport term moves and deforms that field, and the
viscous term acts on its spatial curvature.  All of these terms describe one
motion at the same time.

The same coupling remains when the equation is differentiated.  Let
\(V_n=\partial_t^nu\) and \(P_n=\partial_t^np\).  Differentiating the momentum
equation \(n\) times and applying the product rule to the transport term gives
\[
  V_{n+1}
  +\sum_{k=0}^{n}{n\choose k}(V_k\cdot\nabla)V_{n-k}
  +\nabla P_n
  =\nu\Delta V_n,
  \qquad
  \nabla\cdot V_n=0.
\]
The binomial sum contains every way the \(n\) time derivatives can split across
the two velocity factors in the transport term.  The remaining terms are the
pressure derivative, viscous response, and incompressibility condition at the
same level.  No derivative level evolves on its own.

This is why one property cannot be asked to carry the whole smoothness proof.
The energy identity from Section 1.2 controls the \(L^2\) size of the velocity
and the time-integrated \(L^2\) size of its first derivatives.  It leaves the
pointwise growth of higher derivatives unresolved.  Energy control is real,
but it reaches only part of the structure required for continuation.

Viscosity gives the second example.  The term \(\nu\Delta V_n\) damps spatial
variation at the \(n\)-th time level.  The same equation also contains the full
transport sum and the pressure response.  An estimate that keeps only the
viscous term has dropped other terms acting on the same derivative.  It cannot
describe the evolution of that derivative by itself.

The full VPI law keeps these parts attached to one fluid and one history.  A
Dead, Jump, or Blown endpoint would have to break that joined evolution
somewhere in its derivative structure before it could count as a class exit.
A failure in one selected quantity is only partial evidence.  The Silver route
needs the point where the proposed nonsmooth state can no longer satisfy the
whole law as the same coherent field.

\section{Proof program ((demonstrated and explained with example)(full top-level reasoning and train-of-thought)(reader facing walkthrough))}
